% sec-05: the diligence bank and the failure combinations.  Figure 12
% (graphviz-type fault tree with AND and OR gates), Table 19.
\section{What a Diligence Process Will Ask, and What Would Stop This}

Twenty-two questions, eighteen of which have a file that answers them and four of
which do not. The four are marked rather than distributed among the eighteen,
because a diligence process that finds an unanswered question the company had not
noticed draws a conclusion about the company, and one that finds a question the
company has already routed draws a different one.

\begin{apptable}
\begin{tabularx}{\textwidth}{>{\raggedright\arraybackslash}p{5.4cm} >{\raggedright\arraybackslash}p{3.0cm} >{\raggedright\arraybackslash}X}
\headrow \thead{The question with no answer today} & \thead{Waiting on} & \thead{The route already in motion}\\
\midrule
Which FDA center leads the review of a drug, a robotic platform, and advisory software together & A Pre-Request for Designation response & The letter of this day, held for the next open session \\
Is a perioperative drug supply route available on any path & A developer response to a first approach & The developer letter of day 2, the first approach of any kind \\
Who holds the investigational new drug application & A sponsor-investigator determination & The clinical trial support services letter of day 3 \\
What robotic configuration is proposed & A site agreement & Question 4 of the day 3 feasibility agenda, where the configuration is settled \\
\end{tabularx}
\end{apptable}
\vspace{-0.60cm}
\tabcap{Table 19. The four diligence questions with no answer today, each with\\
what it waits on and the route by which it is already being asked.}

The other eighteen are answered by files in this repository and are listed in
full in \appfile{funding/auto-fund/07Sep26/briefs}, with the file for each named
beside it.

\subsection{The failure combinations}

\begin{appfloat}[!tb]
\begin{appfig}[graphviz-type / fault tree with AND and OR gates]
% Three levels.  The decisive leaf hangs to the left at a shallower depth than
% the AND gates, which keeps its edge from crossing either gate.  Every other
% edge is a straight segment between a gate and a node directly below it.
\node[gvboxk,text width=46mm] (top) at (5.60,0) {The five-day block produces no usable outcome};
\umlgateor{5.60}{-1.05}
\draw[gvundir] (top.south) -- (5.60,-0.79);
\node[gvboxg2,text width=34mm] (dec) at (1.30,-2.62) {No perioperative supply route exists on any path};
\draw[gvundir] (5.28,-1.27) -- node[midway,above,sloped,font=\tiny] {alone is decisive} (dec.north east);
\umlgateand{4.30}{-2.30}
\umlgateand{9.30}{-2.30}
\draw[gvundir] (5.60,-1.27) -- (4.30,-2.04);
\draw[gvundir] (5.92,-1.27) -- (9.30,-2.04);
\node[gvboxs,text width=32mm] (l1) at (3.10,-3.75) {No site will host an investigator-initiated trial};
\node[gvboxs,text width=32mm] (l2) at (5.60,-3.75) {No federal mechanism accepts the shape};
\node[gvboxs,text width=32mm] (l3) at (8.10,-3.75) {A combination determination puts review out of reach};
\node[gvboxs,text width=32mm] (l4) at (10.60,-3.75) {No private instrument closes the gap};
\draw[gvundir] (4.30,-2.52) -- (l1.north);
\draw[gvundir] (4.30,-2.52) -- (l2.north);
\draw[gvundir] (9.30,-2.52) -- (l3.north);
\draw[gvundir] (9.30,-2.52) -- (l4.north);
\node[font=\tiny\itshape,text=fundgray,anchor=north west,text width=118mm]
  at (-0.30,-4.65) {No probability is attached to any leaf. Assigning one would
  invent precision the evidence does not support, and a fault tree is useful for
  its structure whether or not its leaves carry numbers.};
\end{appfig}
\vspace{-0.60cm}
\figcaption{Figure 12. What has to fail for the week to produce nothing, by combination:\\
one decisive single point, and two pairs that each need both halves to fail.}
\end{appfloat}

\subsection{What the tree argues}

That there is exactly one single point of failure, and it is drug supply for the
perioperative investigational use. Everything else requires a pair. That
asymmetry is why the developer letter went out on day 2 rather than being
deferred, and why the Pre-Request for Designation was written on the quietest day
of the block rather than squeezed into a busy one \cite{fdacombination2026}.

The tree also argues something reassuring, which a fault tree usually does not.
Two of the four paired failures already have a second route open: a second
clinical site was approached in parallel on day 3, and a second class of funder,
the disease foundations, was approached on the same day. A pair that requires
both halves to fail is a pair with two chances to survive.

\subsection{Why the stop conditions are published at all}

A funder asked to accept a risk is entitled to know which risks the company has
already identified. A company that names its own stop conditions is easier to
believe about everything else, and the capitalization plan states the same
principle in its own risk section rather than confining it to an appendix
\cite{capitalplan2026}.
